\chapter{Specification of the 3/2 Volatility Model}
\label{chap:model}

\section{Introduction to the 3/2 Model}
\label{sec:model_setup}

As an advanced stochastic volatility framework, the 3/2 model is a complex evolutionary form of the Square-Root Process (popularized by Cox-Ingersoll-Ross and later further developed by Steven Heston).

In the current financial industry, the Heston model (i.e., the 1/2 model) has long been regarded as the industry standard. However, when facing extreme market stress, the Heston model often fails to accurately capture the ``volatility of volatilit'' observed in the market and struggles to perfectly fit the steepness of the short-term volatility skew.

To overcome this limitation, researchers proposed a completely new approach in the late 1990s and early 2000s: modifying the exponent of the diffusion term from 1/2 to 3/2. This specification change allows the model to generate more explosive volatility dynamics and a nonlinear mean reversion speed. Consequently, the 3/2 model allows volatility to spike and presents a steeper skew under market stress, which highly aligns with actual empirical data. However, it should be noted that under certain normal market conditions, the 3/2 model may lack the stability inherent in the Heston model.

\section{Stochastic Differential Equation (SDE) Specification}

Under the risk-neutral measure, the asset price process and the instantaneous variance/volatility process of the 3/2 model can be characterized by the following system of stochastic differential equations:

The dynamics of the underlying asset price process $S_t$ satisfy:

\begin{equation}
    dS_t = r S_t dt + S_t \sqrt{V_t} (\rho dW_t^1 + \sqrt{1-\rho^2} dW_t^2)
    \label{eq:underlyingProcess}
\end{equation}

Where $r$ represents the risk-free interest rate, $W_t^1$ and $W_t^2$ are two independent standard Brownian motions, and $\rho$ is the correlation coefficient between the two Brownian motions. Meanwhile, the dynamic evolution of the variance process $V_t$ satisfies:
\begin{equation}
    dV_t = \kappa V_t (\theta - V_t) dt + \epsilon V_t^{3/2} dW_t^1
    \label{eq:varianceProcess}
\end{equation}

In the above equations:

\begin{itemize}
    \item $\kappa$ controls the speed of mean reversion.
    \item $\theta$ represents the long-term average variance level.
    \item $\epsilon$ is the volatility of volatility parameter, controlling the degree of fluctuation of the variance process.
    \item $W_t^1$ and $W_t^2$ are two independent standard Brownian motions driving the randomness in the variance process and the price process, respectively.
\end{itemize}

It is worth noting that this variance process can also be equivalently expressed in terms of $dV_t/V_t$, thereby more intuitively reflecting the drift and diffusion characteristics of its relative changes.

\section{Integrated Variance and Conditional Distribution}

In the option pricing framework of the 3/2 model, the integrated variance $V_T$ over the time interval $[0, T]$ plays a crucial role. It is defined as follows:

\begin{equation}
    V_T = \int_0^T V_t dt
    \label{integratedVarianceDefinition}
\end{equation}

Since the randomness of the asset price is partially driven by volatility, when we condition on the integrated variance $V_T$, the terminal price of the underlying asset $S_T$ follows a lognormal distribution. This important property means that, given $V_T$, we can borrow the pricing framework of the Black-Scholes (BS) model.

According to Itô's isometry principle:

$$\int_0^T\rho^*\sigma_tdX_t\sim\rho^*\sqrt{V_T}X_1\sim N(0,(\rho^*)^2V_T)$$

Accordingly, the effective volatility $\sigma$ between time $0$ and $T$ can be expressed as:

$$\sigma=\rho^*\sqrt{V_T/T}$$

Through further stochastic calculus derivation, the log-return formula of the terminal price relative to the initial price can be expanded as:
\begin{equation}
    \log\left(\frac{S_T}{S_0}\right) = \frac{\rho}{\epsilon}\log\left(\frac{V_T}{V_0}\right) + \left(\kappa + \frac{\epsilon^2}{2}\right)V_T - \kappa\theta T - \frac{1}{2}V_T + \rho^* \sqrt{V_T} X_1
    \label{priceFormulaForConditionalMC}
\end{equation}

This formula is the core foundation for pricing via subsequent exact simulation methods.

\section{Moment Generating Function of the 3/2 Model}

Under the risk-neutral measure $\mathbb{Q}$, the stochastic differential equations (SDEs) of the 3/2 model are defined by Equation \eqref{eq:underlyingProcess} and Equation \eqref{eq:varianceProcess}.

To utilize Fourier transform methods for option pricing, the core lies in solving the moment generating function (MGF) of the log-asset price $x_T=\ln S_T$.

According to the derivations by Lewis (2000) and Carr and Sun (2007), let $\phi(u,\tau)=\mathbb{E}[e^{u x_T}|\mathcal{F}_t]$ be defined as the moment generating function of the log-price. Within the 3/2 model framework, because the reciprocal of its variance process $V_t$ follows a CIR process (square-root process), this MGF does not possess the common exponential-affine form found in affine models. Instead, it manifests as an analytical form involving Kummer's Confluent Hypergeometric Function ($_1F_1$).

Specifically, given the initial state $V_0=\sigma_0^2$, the auxiliary variables are defined as follows: Parameter transformation terms:
$$\mu=\frac{1}{2}+\frac{\kappa-u\rho\epsilon}{\epsilon^2},\quad\tilde{c}=\frac{u(1-u)}{\epsilon^2}$$

Characteristic parameters:
$$\delta=\sqrt{\mu^2+\tilde{c}},\quad\alpha=-\mu+\delta,\quad\beta=1+2\delta$$

Time-dependent term: define
$$X=\frac{2\kappa\theta}{\epsilon^2\sigma^2(e^{\kappa\theta\tau}-1)}$$

This term characterizes the nonlinear impact of the mean reversion strength and the maturity $\tau$ on the volatility structure. In summary, the moment generating function of the log-price under the 3/2 model can be expressed as:

$$M(u,\tau)=\frac{\Gamma(\beta-\alpha)}{\Gamma(\beta)}X^\alpha\,_1F_1(\alpha,\beta,-X)$$

In the process of numerical computation, the calculation of $_1F_1(\alpha,\beta,-X)$ often involves the processing of special functions in the complex domain. In code implementation, to prevent numerical overflow in large-parameter scenarios (such as extreme volatility or long maturities), the log-gamma function (Log-Gamma) is used for reconstruction. For the MGF pre-factor:
$$M(u,\tau)=\frac{\Gamma(\beta-\alpha)}{\Gamma(\beta)}\,X^\alpha\,_1F_1(\alpha,\beta,-X)$$
Use
$$\ln P=\ln\Gamma(\beta-\alpha)-\ln\Gamma(\beta)+\alpha\ln X,\qquad P=\exp(\ln P),$$
then
$$M=P\cdot\,_1F_1(\alpha,\beta,-X)$$
to replace direct multiplication/exponentiation. This approach guarantees the numerical stability of the operations.